\section{État de l'art : De la prédiction de mots à l'Agentification}

\subsection{Émergence des LLM}
Les premières idées de réseaux de neurones datent des années 40, mais ceux-ci ne connaissent pas un franc succès. À partir des années 80, le développement de l’algorithme de rétropropagation leur permet de revenir sur le devant de la scène mais toujours sans grandes réussites car les modèles ont alors une taille intermédiaire qui limite les possibilités de cette technologie. Ce n’est qu’au début des années 2010 que les performances des réseaux de neurones ne cessent d’augmenter grâce à l’amélioration des capacités de calcul des GPU qui permet l’émergence des réseaux de neurones profonds et la mise en place de nouvelles méthodes.
\\

Si des méthodes de Natural Language Processing (NLP) à l’aide de réseaux de neurones sont étudiées depuis longtemps, le traitement et la génération de texte ont connu des progrès phénoménaux depuis 2017 avec l’implémentation du mécanisme d’attention et des transformers, dans l’article \textit{Attention is all you need} \cite{Vaswani2017}, menant à la sortie de ChatGPT 3 en 2022. Ces progrès se sont poursuivis avec l’apparition de modèles de raisonnements, toujours plus efficaces notamment pour la génération de code.
\\

Afin d’améliorer les performances des LLM, la tendance actuelle semble être à l'"agentification". Une IA agentique améliore itérativement la réponse d’un LLM en lui donnant accès à une grande variété d’outils (ex: une calculatrice, un terminal pour tester le code avant de le fournir à l’utilisateur, la recherche internet, etc) et en lui permettant d’adapter progressivement sa réponse grâce aux résultats obtenus à partir des outils. La sortie de ChatGPT-5 est un grand pas dans cette direction car c’est un agent IA qui rassemble en un seul modèle tous les modèles de la génération précédente d'OpenAI en choisissant de façon autonome lequel est le plus adapté à la requête de l’utilisateur.

\subsection{Hallucinations et leurs solutions}

Cependant, les résultats des LLM sont parfois limités : les LLM étant entraînés pour donner une réponse quoi qu’il arrive, ils ont tendance à proposer une réponse fausse avec certitude face à des problèmes difficiles, ce que l’on appelle des hallucinations \cite{Tauman2025}. C’est d’autant plus le cas pour des DSL, que les LLM connaissent mal, et si la seule documentation du langage suffit à leur faire produire du code syntaxiquement correct, celui-ci donne rarement le bon résultat.
\\

Une première solution simple à implémenter pour éviter les hallucinations est de formuler la requête donnée au LLM avec précision et ingéniosité. Il a par exemple été remarqué qu’on obtient de meilleurs résultats en moyenne si on précise au modèle de langage son rôle; c’est ce qu’on appelle le \textit{Persona Pattern}. Ainsi la requête: “Tu es mon professeur de mathématiques, explique moi et démontre moi l’inégalité de Cauchy-Schwarz” donnera probablement un meilleur résultat que la simple requête : “Explique-moi et démontre-moi l’inégalité de Cauchy-Schwarz”. Cette astuce est une des nombreuses techniques qui constituent le \textit{Prompt Engineering} \cite{White2023}.
\\

Avec l’émergence des agents IA évoquée plus haut, une nouvelle couche de complexité s’est ajoutée au Prompt Engineering, on parle alors de \textit{Context Engineering}. Cette discipline ne se contente plus d'optimiser la formulation des instructions, mais se focalise sur la sélection et l'organisation stratégique des données fournies au modèle. L'enjeu principal réside dans la gestion du rapport signal-sur-bruit : un excès d'informations peut "distraire" l'agent et dégrader ses performances, tandis qu'un contexte trop restreint limite sa capacité d'action. 

Le Context Engineering repose ainsi sur une structuration rigoureuse (notamment via l'usage du format JSON ou encore de balises XML) pour distinguer les outils, les documents et les consignes, ainsi que sur des mécanismes de filtrage dynamique. En optimisant ce flux informationnel, on permet à l'agent de naviguer efficacement dans des workflows itératifs complexes, tout en maîtrisant la latence et les coûts grâce à des techniques comme la mise en cache du contexte \cite{Anthropic2025}. 
\\

Une approche alternative, plus fondamentale mais néanmoins digne d'intérêt, consisterait à utiliser des outils de recherche lexicale, dont la commande grep est l'exemple le plus emblématique. Cette méthode se base sur la recherche de correspondances exactes de chaînes de caractères ou d'expressions régulières directement dans les fichiers de code source. Si cette technique est dépourvue de toute compréhension du sens — elle ne ferait pas le lien entre "calculer le stock" et une variable nommée "inventory-level" — elle offre une rapidité et une précision inégalées pour des requêtes littérales, comme retrouver toutes les occurrences d'un nom de fonction ou d'une variable spécifique.
\\

Une autre solution possible à ce problème est d’utiliser des méthodes de \textit{Retrieval Augmented Generation} (RAG) \cite{Gao2023}, qui consistent à retrouver les éléments de documentation liés à la requête de l’utilisateur et les fournir directement au LLM pour améliorer ses réponses. Pour cela, on peut par exemple utiliser l’indexation sémantique avancée, en représentant le texte sous forme de vecteurs (appelés \textit{embeddings}) tels que deux textes sémantiquement similaires soient représentés par des vecteurs proches \cite{Reimers2019}. Des modules comme \textit{Sentence Transformers} permettent de générer ces vecteurs, et d’autres tels que \textit{FAISS} (Facebook AI Similarity Search) servent à formuler des recherches de similarité sur des vecteurs de grande dimension via des calculs de proximité selon différentes métriques.

Cependant, la recherche sémantique classique par similarité cosinus présente des limites face à des requêtes asymétriques ou du vocabulaire très spécifique. Pour pallier ce fossé sémantique, des techniques d'augmentation de requêtes telles que HyDE (Hypothetical Document Embeddings) ont vu le jour \cite{Gao2022}. Cette méthode exploite un LLM pour générer une réponse hypothétique à la question de l'utilisateur, et c'est le vecteur de cette réponse générée qui est ensuite utilisé pour chercher les documents pertinents dans la base. Le modèle d'encodage agit alors comme un filtre qui élimine les informations factuellement fausses générées dans l'hypothèse, tout en capturant les motifs sémantiques pertinents.

Dans le contexte de bases de code complexes, il est souvent nécessaire de comprendre la structure globale d'un projet plutôt que de lire de simples fragments de code isolés. L'architecture RAPTOR (Recursive Abstractive Processing for Tree-Organized Retrieval) répond à ce besoin en regroupant et en résumant récursivement les fragments de texte pour construire un arbre hiérarchique de connaissances \cite{Sarthi2024}. Lors de la recherche, le système peut ainsi extraire des informations à différents niveaux d'abstraction, ce qui améliore significativement les performances sur des questions nécessitant une compréhension holistique et un raisonnement multi-étapes.

Enfin, pour maximiser la précision contextuelle, la simple récupération de documents initiaux est aujourd'hui supplantée par des architectures intégrant un reclassement (reranking) neuronal. L'utilisation de modèles d'attention croisée (Cross-Encoder), popularisée par les travaux de Nogueira et Cho \cite{Nogueira2019}, permet d'évaluer la pertinence en traitant simultanément la question et le document candidat au sein du même réseau. Bien que plus coûteuse en calcul, cette méthode capture des relations sémantiques fines inaccessibles aux simples recherches vectorielles et s'avère indispensable pour hisser les extraits de code les plus exacts en tête de liste.

Peu importe la technique utilisée, les données ainsi extraites devant être fournies directement dans la requête, la formulation précise de cette dernière est un aspect central du RAG \cite{Rice2025}. Le RAG utilise donc plusieurs techniques de Prompt Engineering.
\\

Une dernière possibilité est le fine-tuning de modèles pré-entraînés : plutôt que de réentraîner un modèle spécialisé à partir de rien, ce qui est très coûteux, il est possible de spécialiser l’entraînement d’un modèle généraliste pour l’adapter à un cas d’usage spécifique \cite{Pratap2025}. Cependant, pour mettre en place cette solution efficacement, il faut une grande quantité d’exemples de questions/réponses. Cela peut donc conduire à l’utilisation d’un autre LLM (souvent plus gros) afin de générer d’autres exemples \cite{Mulder2023}. Pour ces raisons, par manque de temps et afin de préserver la confidentialité des données de Lokad, nous n'avons pas implémenté de fine tuning dans notre PSC.